\inicializarSeccion{Gráficos relevantes Encontrados}
Este gráfico nos mostró que hay un rango de medidas deseable para nuestro producto, por lo que impactó en nuestra decisión final sobre las medidas, tomando el rango de 8.50 cm - 10.50 cm como ganador.

\begin{figure}[H]
    \centering
    \includegraphics[width=0.8\textwidth]{\nat{parte2/comparacionDeDiametros.png}}
    \caption{Gráfico de medidas del diámetro (Pregunta 9).}
\end{figure}

Otro gráfico muy relevante fue el gráfico de pay de la pregunta seis. 

\begin{figure}[H]
    \centering
    \includegraphics[width=0.5\textwidth]{\nat{parte2/analisisPregunta6.png}}
    \caption{Gráfico de pastel que muestran las opciones de material; la moda fue acero inoxidable.}
    \label{fig:analisisPregunta6}
\end{figure}

Para este gráfico nos pareció muy impactante la preferencia por el material de acero inoxidable, por eso nuestra decisión de usar este material sobre cualquier otro, independientemente del costo que tenga. 


Por último, el gráfico de barras de la pregunta 2 también nos pareció muy relevante, ya que le preguntamos al cliente de manera indirecta sobre el tamaño. Si es que quiere algo mas elegante y compacto, o quizas algo mas robusto pero que pueda conservar mejor la temperatura.

\begin{figure}[H]
    \centering
    \includegraphics[width=0.8\textwidth]{\nat{parte2/visualizacionDelAnalisis.png}}
    \caption{Gráfico de barras que muestran las opciones de aspectos para el diseño (Pregunta 2).}  
\end{figure}